\section{Thinning by a subset of the alphabet}
\label{sec:poisson_word_partition}

% Derived supporting identities for the label selection in
% 09-amplification.tex at adc7f1241b42e322a6451854ab7e4b4c146bf78a,
% lines 49--54 and 237--253; not separately stated theorems in that source.
The retained and omitted label words associated with a subset of the
alphabet satisfy the following fixed-time identities. The retained labels
in the local constraint products of \cite{OpenAI2026AreaLaw}, in the section
``Amplifying the collar estimate,'' are an example of such a subset.

\begin{definition}[Partitioning a word]
\label{def:poisson_word_partition}
\lean{PoissonWord.partitionWords, PoissonWord.partitionWords_nil,
PoissonWord.partitionWords_length}
\leanok
\uses{def:poisson_word_relabel, def:poisson_word_filters}
Let $p$ be a predicate on an alphabet $I$, and put
$S=\{i\in I:p(i)\}$ and $T=\{i\in I:\neg p(i)\}$.
Classify each letter by the bijection $c:I\to S\sqcup T$, then define
$P_p(w)=(L(c^{\mathcal W}(w)),R(c^{\mathcal W}(w)))$.
Write $P_p(w)=(w_S,w_T)$. The empty word gives two empty words, and
$|w_S|+|w_T|=|w|$.
\end{definition}

\begin{theorem}[The original filtered letters]
\label{thm:poisson_partition_letters}
\lean{PoissonWord.ofFn_partitionWords_fst, PoissonWord.ofFn_partitionWords_snd}
\leanok
\uses{def:poisson_word_partition}
For $w=(i_0,\ldots,i_{m-1})$, regard the letters of $w_S,w_T$ as letters
of $I$ through their inclusions. Then
\begin{align}
 (w_S,w_T)&=\bigl((i_j)_{0\leq j<m,\ p(i_j)},
                  (i_j)_{0\leq j<m,\ \neg p(i_j)}\bigr),
 \label{eq:poisson_partition_letters}
\end{align}
where both lists retain increasing order of the original positions.
\end{theorem}
\begin{proof}
\leanok
\uses{def:poisson_word_partition, def:poisson_word_filters,
def:poisson_word_relabel}
For each letter $i$, the left filter retains $c(i)$ exactly when $p(i)$
holds, and the right filter retains it exactly when $\neg p(i)$ holds.
In either case the inclusion into $I$ returns $i$. Applying this identity
at every position proves both equalities, preserving the original order.
\end{proof}

\begin{theorem}[Independent subset thinning]
\label{thm:poisson_word_partition}
\lean{PoissonWord.map_partitionWords, PoissonWord.map_partitionWords_fst,
PoissonWord.map_partitionWords_snd, PoissonWord.indepFun_partitionWords}
\leanok
\uses{def:poisson_word_partition}
Suppose $I$ is finite and $t\geq0$. For a single sample $w$ with law
$\mu_{I,t}$, its two ordered subwords have joint law
\begin{align}
 (P_p)_*\mu_{I,t}&=\mu_{S,t}\otimes\mu_{T,t}.
 \label{eq:poisson_partition_law}
\end{align}
In particular, $w_S$ and $w_T$ are independent with laws $\mu_{S,t}$ and
$\mu_{T,t}$. This holds when either subset is empty and at $t=0$.
\end{theorem}
\begin{proof}
\leanok
\uses{thm:poisson_word_relabel, thm:poisson_word_thinning,
lem:poisson_word_normalization}
Relabeling gives $(c^{\mathcal W})_*\mu_{I,t}=\mu_{S\sqcup T,t}$.
Applying the joint filtered-word law to this measure gives
\begin{align}
 (P_p)_*\mu_{I,t}
 &=(L,R)_*(c^{\mathcal W})_*\mu_{I,t}
  =(L,R)_*\mu_{S\sqcup T,t}
  =\mu_{S,t}\otimes\mu_{T,t}.
\end{align}
Both factors have mass one, so coordinate projections give the stated
marginals. Their product is the joint law, which proves independence.
\end{proof}
